\documentclass[12pt]{article}
\usepackage[utf8]{inputenc}
\usepackage[T1]{fontenc}
\usepackage[margin=1in]{geometry}
\usepackage{setspace}
\usepackage{xcolor}
\usepackage[colorlinks=true, urlcolor=blue!60!black, linkcolor=blue!60!black]{hyperref}
\setlength{\parindent}{2em}
\setlength{\parskip}{1.25em}
\renewcommand{\baselinestretch}{1.0}
\begin{document}
\begin{center}
{\LARGE \textit{Puerto Rico v. Branstad}, 483 U.S. 219 (1987)\par}
\vspace{0.8em}
{\large Case Brief --- Four-Page Edition\par}
\end{center}
\vspace{1em}

\section{Memory Jogger}

When Iowa's Governor refused to extradite a man charged with murder in Puerto Rico, the Court held unanimously that Puerto Rico could obtain a federal court order enforcing the Governor's mandatory duty under the Extradition Act, overruling \textit{Kentucky v.\ Dennison} to the extent it made that statutory duty unenforceable. Six Justices reached the same conclusion under the Extradition Clause (\href{https://supreme.justia.com/cases/federal/us/483/219/}{483 U.S. 219}).

\section{Detailed Case Facts}

In January 1981, Ronald Calder, an Iowa native working as a federal air traffic controller in Puerto Rico, struck two pedestrians with his car in Aguadilla after an altercation with Antonio de Jesus Gonzalez. De Jesus Gonzalez was injured, and his wife, Army Villalba, eight months pregnant, was killed with her unborn child; two sworn statements said Calder backed his car over her body two or three times. Charged with first-degree murder and attempted murder, Calder twice failed to appear, was declared a fugitive, and surrendered in Iowa, where he posted bond.

Puerto Rico's Governor sent a formal requisition. At an extradition hearing, Calder's attorney testified that ``a white American man . . . could not receive a fair trial in the Commonwealth of Puerto Rico.'' Governor Robert Ray denied extradition, citing the absence of a ``change to a more realistic charge,'' and his successor, Governor Terry Branstad, also refused.

\section{Procedural History}

After the Supreme Court denied leave to file an original action, Puerto Rico sued Branstad and Iowa in the Southern District of Iowa for mandamus under the Extradition Clause and 18 U.S.C. \S~3182. Although Iowa stipulated that the papers satisfied the statute, the district court dismissed under \textit{Kentucky v. Dennison} (1861). The Eighth Circuit ``[r]eluctantly'' affirmed, and the Supreme Court reversed on June 23, 1987.

\section{Judicial Votes}

Reversed, \textbf{9--0} on the judgment (6--3 on the constitutional holding of Part II-B).
\begin{itemize}
\item \textbf{Majority:} Justice Marshall, joined in full by Chief Justice Rehnquist and Justices Brennan, White, Blackmun, and Stevens; by Justices Powell and O'Connor except Part II-B; and by Justice Scalia as to the Extradition Act and the overruling of \textit{Dennison}'s reading of its predecessor.
\item \textbf{Concurring in part and in the judgment:} Justice O'Connor, joined by Justice Powell; Justice Scalia.
\item \textbf{Dissents:} None.
\end{itemize}

\section{Holding}

Overruling \textit{Dennison}, federal courts may issue mandamus compelling an asylum State's governor to perform the mandatory, ministerial duty to deliver up a fugitive upon proper demand under the Extradition Clause and the Extradition Act. Without deciding whether the Clause itself covers Puerto Rico, the Court held that \S~3182 applies to Puerto Rico as a Territory, so it may seek federal mandamus against a state governor.

\section{Analysis of Opinions}

\textbf{Majority (Marshall): Dennison in context.} Decided on March 14, 1861, after secession and weeks before Fort Sumter, \textit{Dennison} arose when ``the practical power of the Federal Government [was] at its lowest ebb since the adoption of the Constitution.'' Taney held that Ohio's Governor had an absolute, ``merely ministerial'' duty to surrender a free Black man indicted for helping an enslaved woman escape, yet called the duty merely ``moral,'' because ``the Federal Government, under the Constitution, has no power to impose on a State officer, as such, any duty whatever, and compel him to perform it.''

\textbf{Obsolete federalism.} Marshall reaffirmed that the Clause is mandatory, consistent with \textit{Michigan v. Doran}, but repudiated Taney's premise of coequal sovereigns. Federal courts routinely enforce federal duties against state officers by mandamus and injunction (\textit{Ex parte Young}, \textit{Brown}, \textit{Cooper v. Aaron}). Because the duty is ministerial, explicit, and constitutionally imposed, enforcing it requires no ongoing supervision and no Tenth Amendment balancing, and entrenched noncompliance cannot create an ``executive common law'' of discretion.

\textbf{Puerto Rico's statutory right.} Section 3182 covers any demanding ``Territory''; under \textit{Kopel v. Bingham}, Puerto Rico qualified and did not lose that authority as a Commonwealth, and federal courts may enforce federal statutory duties against state officials (\textit{Maine v. Thiboutot}).

\textbf{O'Connor (concurring in part).} Joined by Powell, she invoked constitutional avoidance: the statute resolved the case, so Part II-B's pronouncements were unnecessary.

\textbf{Scalia (concurring in part).} He joined only the statutory holding and the overruling of \textit{Dennison}'s reading of the 1793 Act, noting ``that no party before us has asserted the lack of power of Congress to require extradition from a State to a Territory.'' \textit{New Mexico ex rel. Ortiz v. Reed} (1998) later applied \textit{Branstad} to forbid refusals based on fear of unfair treatment.

\section{Examples: Future Applications}

\begin{enumerate}
\item \textbf{Same-Side (Compelled): Shield Law.} A doctor who performed a banned procedure while physically in State X flees to State Y, whose shield law bars extradition. The Clause and \S~3182 preempt it.
\item \textbf{Same-Side (Compelled): Death-Penalty Objection.} State Y's Governor demands that State X forgo a death sentence. Like Iowa's fair-trial doubts, a moral objection is no defense under \textit{Doran}.
\item \textbf{Opposite-Side (No Duty): Telehealth Prescriber.} A doctor who never entered State X never ``fled from Justice,'' and a Commerce Clause commandeering statute would likely fall under \textit{Printz} and \textit{Murphy}.
\item \textbf{Opposite-Side (No Duty): Foreign Demand.} A Canadian province's request runs through treaties, not Article IV or \S~3182.
\item \textbf{Fence-Sitter: Puerto Rico After \textit{Printz}.} A duty to surrender fugitives to a Territory rests on statute alone, raising the commandeering question Scalia reserved.
\end{enumerate}

\section{Critique}

\textbf{Analytical mismatch.} As O'Connor and Scalia flagged, Part II-B relies on a duty ``directly imposed upon the States by the Constitution itself'' to bypass the Tenth Amendment, yet Part II-C rests Puerto Rico's remedy on a statute alone. If extension to Territories is purely statutory, \textit{Printz}'s anti-commandeering rule creates genuine tension with \textit{Branstad}'s Territory holding.

\textbf{Horizontal federalism.} Mandatory rendition promotes interstate harmony and prevents sanctuaries, yet strips asylum governors of any humanitarian safety valve when demanding states pursue contested prosecutions, the dilemma that animated antebellum free-state governors in \textit{Dennison}.

\section{Key Quotations}

\begin{quote}
``\emph{Kentucky v. Dennison} is the product of another time. The conception of the relation between the States and the Federal Government there announced is fundamentally incompatible with more than a century of constitutional development. Yet this decision has stood while the world of which it was a part has passed away. We conclude that it may stand no longer.''
\end{quote}
\noindent --- Justice Marshall, Opinion of the Court, 483 U.S. at 230.

\begin{quote}
``We reaffirm the conclusion that the commands of the Extradition Clause are mandatory, and afford no discretion to the executive officers or courts of the asylum State.''
\end{quote}
\noindent --- Justice Marshall, Opinion of the Court, 483 U.S. at 227.

\begin{quote}
``Because the duty is directly imposed upon the States by the Constitution itself, there can be no need to weigh the performance of the federal obligation against the powers reserved to the States under the Tenth Amendment.''
\end{quote}
\noindent --- Justice Marshall, Opinion of the Court, 483 U.S. at 228.

\end{document}
